\begin{pastexam}[label={px:nccu-113-dsa-08}]{政大資科所 113 學年度 第 8 題}
（單選；不必說明理由）

(4\%) In the context of computational complexity, which of the following best describes a
pseudo-polynomial time algorithm?
\begin{enumerate}[label=\Alph*.]
\item An algorithm whose running time is polynomial in the size of the input.
\item An algorithm whose running time is polynomial in the numeric value of the input, often involving
dynamic programming.
\item An algorithm that runs in logarithmic time.
\item An algorithm with linear time complexity, solving problems efficiently regardless of the input size.
\end{enumerate}
\end{pastexam}

\begin{pastsoln}
答案是 B。 資結書定義~28.5：擬多項式時間是輸入長度和輸入中最大的\textbf{數值}兩者的多項式。 B 說的「是輸入
數值的多項式，常用動態規劃」最接近，例如 0/1 背包的 $\Theta(nW)$、子集合和的 $\Theta(mt)$，時間和 $W$、$t$ 的
數值成正比（資結書第~18.5.5、28.6~節）。 B 的說法少了「和輸入長度」，但四個選項裡只有它抓到「數值」這個重點。

A 是多項式時間的定義（資結書第~28.1~節）：時間是輸入大小（位元數）的多項式。 擬多項式時間不一定是多項式時間，
因為數值 $M$ 只占 $\lg M$ 個位元，$\Theta(nM)$ 是位元數的指數。 C、D 只說時間是對數或線性，和「依數值還是依
位元數量」無關；D 的「不管輸入大小都很快」也不對，線性時間仍隨輸入大小增加。
\end{pastsoln}
